% ── 다리 초안: weppner_huggins1977 → 우리 eq:vn (§1 측정원리 bgbox) ──
% 삽입 위치: ch1_sec01_n0n1.tex, 측정원리 bgbox(현행 204–218행) 직후,
%   자산 주석(현행 219행) 앞. (baek_pilon 다리와 인접 배치 — br_baek_pilon2022.tex 참조.)
% 앵커 문장: "...GITT(...)\cite{weppner_huggins1977} --- 이 준평형 OCV 곡선의 정통 측정이고,
%   저율 정전류 곡선의 분극 보정(식~\eqref{eq:vn})은 같은 목표의 연속 근사판이다."(현행 208–209행)
% 컨테이너: srcbox. 우리 기호로 서술.
% ★확인 필요: 원 논문(J. Electrochem. Soc. 124,1569)의 확산계수 √t 식(정확형·식번호) 미확인
%   (PDF 파싱 불가) — 단 본 다리는 D 식이 아니라 $U_\oc$ 측정 대응이 핵이라 D 식 불요.
%   방법 요지(펄스→완화 준평형 OCV·√t 의존)는 검증됨.

\begin{srcbox}
\textbf{다리 --- eq:vn 이 GITT 와 무엇을 공유하나.} 본문이 입력으로 전제하는 준평형 $U_\oc(\bar x)$ 는
Weppner--Huggins 의 GITT\cite{weppner_huggins1977} 가 재는 그 양이고, 식~\eqref{eq:vn} 의 분극 보정
$V_n=V_\app-\sigma_d|I|R_n$ 은 같은 $U_\oc$ 로 가는 연속 근사판이다:
\begin{center}\footnotesize
\begin{tabular}{@{}ll@{}}
\hline
본문 & Weppner--Huggins\cite{weppner_huggins1977}(GITT) 대응 \\
\hline
목표량 준평형 $U_\oc(\bar x)$ & 펄스 후 완화된 정지 전위(준평형 OCV) \\
분극 보정항 $\sigma_d|I|R_n$ (식~\eqref{eq:vn}) & 대기 완화로 제거되는 IR$\cdot$과전압 \\
저율 연속 곡선(연속 근사) & 펄스--완화 이산 점열(정통 측정) \\
\hline
\end{tabular}
\end{center}
\emph{방법 요지} --- Weppner--Huggins 는 정전류 펄스로 조성을 한 걸음 옮긴 뒤 충분히 완화된 정지 전위를 읽어
준평형 OCV 점열을 얻고(펄스 중 전위의 $\sqrt t$ 의존에서 화학 확산계수까지 함께 추출), 그 점열이 준평형
$U_\oc(\bar x)$ 의 정통 측정이 된다. \emph{가정 차} --- GITT 는 각 점에서 완전 완화를 기다려 분극을 $0$ 으로
만드나, 식~\eqref{eq:vn} 은 저율 연속 곡선에서 분극을 상수-저항 $R_n$ 모형($\sigma_d|I|R_n$)으로 근사 제거하므로
--- $R_n$ 의 조성$\cdot$전류 의존을 무시하는 만큼 --- GITT 완화만큼 완전하지 않은 연속 근사다.
\end{srcbox}
